\section{引言}
\label{sec:introduction}

大语言模型已经成为函数级代码生成系统的核心组件。HumanEval 与 MBPP 等基准表明，强模型能够在相当比例的短函数任务上通过单次生成得到可运行实现\cite{chen2021evaluating,austin2021program}。真实部署进一步关心接口计费、上下文长度与调用预算共同约束下的系统编排。对同一失败样本，重新采样、执行反馈修补、结构化重链、候选扩展与执行筛选都会改变通过率和成本之间的关系\cite{li2022alphacode,chen2024selfdebug,shinn2023reflexion,gou2024critic,li2025sstar}。

现有研究已经分别验证了若干有效组件：执行反馈能够修补错误程序，候选扩展与筛选能够提高通过率，反思或批判式提示能够改善复杂任务。然而这些组件在工程系统中往往以多个开关叠加出现，若执行器、统计口径或默认提示模板不同，最终收益很难归因到预算组织本身。本文关注的正是这一层问题：在同一实现中逐级解锁不同测试时计算投向，同时记录通过率、调用次数与令牌成本。

本文将这一口径称为\textbf{预算分档编排}。协议把函数级代码生成流程组织为四个语义明确的档位：单步档作为最低成本基线；单步+修补档在失败后引入执行反馈；分流+重链与救援档依据结构难度先验选择重路径，并将失败信息压缩回注；级联档在高预算下扩展候选集并用执行结果筛选。四档共享同一外层闭环，差异只在新增预算投向，使成本--收益关系能够被直接比较。

该协议同时给出机制层面的解释路径。结构难度分承担预算排序作用，用于把更复杂的样本优先送入重路径；固定轻/重路径多重集的随机对照用于隔离覆盖率和结构对齐；压缩回注与完整回溯的对照刻画反馈形式对修补收益和令牌成本的影响；级联消融进一步区分候选扩展和后处理精修的贡献。这样，端到端结果和模块机制可以在同一实验口径下相互印证。

本文的主要贡献如下：

1）\textbf{协议贡献}：提出面向函数级代码生成的预算分档编排协议，在统一执行器、统一判定与统一记账下比较单步、修补、分流救援与级联四类预算形态，将测试时计算从调用次数多少转化为预算投向差异。

2）\textbf{机制贡献}：围绕分流、反馈与级联三个环节给出机制实验。结构难度分被建模为预算分配先验，固定轻/重路径多重集对照检验结构对齐效果，压缩反馈与完整回溯的五随机种子对照评估回注形式。

3）\textbf{评测贡献}：在 HumanEval 与 MBPP 双基准上报告四档三随机种子主结果，并同时给出通过率、平均调用次数、总令牌与经验收益率 $\rho$。实验显示，MBPP 上分档收益更充分，HumanEval 上主要体现为高位残差收束，二者共同展示了预算编排在不同任务分布下的价值。

全文组织如下。第~\ref{sec:related} 节梳理代码生成评测、执行反馈修补与测试时扩展相关工作；第~\ref{sec:formulation} 节给出问题定义、预算档位和成本记账口径；第~\ref{sec:method} 节描述分流、救援与级联的实现机制；第~\ref{sec:experiments} 节报告主结果与消融；第~\ref{sec:discussion} 节讨论部署含义、机制边界与局限；第~\ref{sec:conclusion} 节总结全文。
